\section{A linear learner-depth cost is necessary}
\label{sec:linear-depth-lower}
A bounded-output Gaussian likelihood cannot quickly undo exponential shrinkage of prior function differences. This gives a linear-depth lower bound for the sparse priors at a fixed sufficiently small variance, even with native biases at every layer. It is a fixed-temperature obstruction, not an all-temperature impossibility theorem.

\begin{lemma}[From shrinking increments to a sample lower bound]
\label{lem:increment-learning-lower}
Let a prior on functions $f_w:\mathcal X\to[-1,1]$ satisfy
\[
 \E_\pi|f_w(x_+)-f_w(x_-)|^2\le K r^H,
 \qquad K>0,\quad0<r<1.
\]
Let $X$ be uniform on the two fixed points $x_+,x_-$ and set
$Y=a$ at $x_+$ and $Y=-a$ at $x_-$, where $0<a\le1$.
Use the ordinary posterior with likelihood
$\exp[-\sum_i(f_w(X_i)-Y_i)^2/T]$, $T>0$.
For every realized sample,
\[
 |\bar f_S(x_+)-\bar f_S(x_-)|\le
       e^{4n/T}\sqrt K\,r^{H/2}.
\]
Whenever the right side is at most $a$, one has $R(\bar f_S)\ge a^2/4$.
In particular, for all sufficiently large $H$, depending only on $K,r,a$, this holds if
\[
 n\le\frac{TH\log(1/r)}{16}.
\]
\end{lemma}
\begin{proof}
Every squared training loss is at most four, so the normalizer is at least $e^{-4n/T}$ and the likelihood is at most one. Apply this density bound and Cauchy--Schwarz to the prior increment. If the two posterior-mean values are $u,v$, their risk is
\[
 \tfrac12[(u-a)^2+(v+a)^2]
 =\left(\frac{u+v}{2}\right)^2+
       \left(a-\frac{u-v}{2}\right)^2.
\]
This is at least $a^2/4$ when $|u-v|\le a$. The specified sample count leaves the increment bound at most $\sqrt K r^{H/4}$, which is eventually at most $a$.
\end{proof}

\begin{proposition}[Independent small-variance weights]
\label{prop:linear-depth-lower}
Consider a width-$N$ network with $H\ge1$ hidden layers, a globally $A$-Lipschitz activation, a linear or affine raw readout, and a final one-Lipschitz link into $[-1,1]$. Weight coordinates are independent and centered. A weight in a row of fan-in $F$ has variance at most $v/F$, where $F$ is at least the number of nonbias inputs. Biases may be present at every layer and may have any fixed independent law. If
\[
 A^2v<1,
\]
then, for $x_\pm=\pm e_1\in[-1,1]^d$,
\[
 \E|f_w(x_+)-f_w(x_-)|^2\le4v(A^2v)^H.
\]
Consequently the noiseless two-input problem of Lemma~\ref{lem:increment-learning-lower} has an $\Omega(TH)$ sample lower bound, uniform in every finite learner width. Neither its input distribution nor its two target values depend on $H,n,$ or $N$.

For ReLU and symmetric weights and biases, the stronger bound is
\[
 \E|f_w(x_+)-f_w(x_-)|^2\le4v(v/2)^H,
\]
so only $v<2$ is needed. In particular, at $v=1/8$, $a=1/2$, every $H\ge1$ and every
$n\le TH\log16/16$ give $R(\bar f_S)\ge1/16$.
With a final clip, the target is the width-two, one-hidden-layer bounded-coefficient ReLU function
$\tfrac12[\operatorname{ReLU}(x_1)-\operatorname{ReLU}(-x_1)]$.
\end{proposition}
\begin{proof}
Biases cancel from preactivation differences. Conditional on the previous layer, centered independent incoming weights eliminate cross terms. Thus if the largest preceding mean-square feature difference is $D_{\ell-1}$, the next preactivation difference has second moment at most $vD_{\ell-1}$, and Lipschitzness gives
$D_\ell\le A^2vD_{\ell-1}$.
At the first layer the two inputs differ in one coordinate by two, giving $D_1\le4A^2v$. The readout contributes one further factor $v$; a final one-Lipschitz link does not increase the difference. This proves the first estimate without any symmetry requirement or zero-at-zero activation assumption.

For ReLU, condition on the previous features. The pair of preactivations $(U,V)$ is jointly centrally symmetric. Since
$\operatorname{ReLU}(u)=(u+|u|)/2$, central symmetry cancels the cross term and gives
\[
 \E[\operatorname{ReLU}(U)-\operatorname{ReLU}(V)]^2
 =\tfrac14\E[(U-V)^2+(|U|-|V|)^2]
 \le\tfrac12\E(U-V)^2.
\]
This remains valid whenever the increment is square integrable, because the absolute-value increment is bounded by $|U-V|$. The same layer recursion now uses $v/2$. With $v=1/8$ and the stated sample count, Lemma~\ref{lem:increment-learning-lower} bounds the posterior-mean difference by
$2\sqrt{1/8}\,16^{-H/4}=2^{-H}/\sqrt2<1/2$.
\end{proof}

The L\'evy fan-in prior with intensity
$\mathfrak n(du)=(1/8)|u|^{-1}(1+|u|)^{-3}du$
has variance exactly $1/(8F)$ in a time-$1/F$ coordinate, because
$2\int_0^\infty u(1+u)^{-3}du=1$.
It therefore satisfies this lower bound. Signed beta-prime fan-in laws satisfy it too when their fixed shape constant is small enough. A bound on first absolute moments alone does not imply the required variance bound; the second moment must be checked for the chosen scalar family.

\paragraph{The neuron-gated prior.}
The same criterion applies to the hierarchical gamma-gate model in its contractive variance regime. Let $\mu_2=\E G^2$ for its centered independent edge factors, and use $S_j\sim\operatorname{Gamma}(\kappa/N,1)$. Own-layer gates are independent of their incoming features, and centered factors eliminate cross terms, so the internal mean-square gain is at most
\[
 v_*=N\E S_j^2\,\mu_2
      =\kappa(1+\kappa/N)\mu_2
      \le\kappa(1+\kappa)\mu_2.
\]
The first-layer feature difference has second moment at most $4A^2\mu_2$; the raw readout adds a factor $v_*$. Hence the bounded output increment is at most
$4A^2\mu_2v_*(A^2v_*)^{H-1}$ for a one-Lipschitz output link. A fixed sufficiently small $\kappa$ makes $A^2v_*<1$, uniformly in $N$, and gives the same linear-depth lower bound. With a fixed output-link Lipschitz constant, multiply the displayed bound by its square. This statement concerns that variance regime; it is not asserted for every gate/factor parameter choice.

\paragraph{A raw-likelihood version without native internal biases.}
For ReLU with symmetric independent small-variance coordinates, no native internal or readout bias, and a constant input in the first layer (counted in its fan-in, so $F\ge d+1$), one also has
$\E z_w(x)^2\le v(v/2)^H$ uniformly on the cube.
Take the constant noiseless target $Y=1/2$. For the raw Gaussian likelihood, Jensen gives
\[
 Z_S\ge\exp\{-n[1/4+v(v/2)^H]/T\}.
\]
Thus the posterior mean of $|z_w(x)|$ is bounded by this inverse normalizer times $\sqrt v(v/2)^{H/2}$, and again constant error persists for $n\le cTH$ at fixed $v<2$. The same bound applies to posterior means of individually clipped draws and to clipping the raw posterior mean. The absence of native internal/readout biases is needed for this absolute-output decay argument, whereas the bounded-likelihood increment result above allows biases.

\begin{proposition}[Raw likelihood and prediction, including native biases]
\label{prop:raw-increment-depth-lower}
Let a prior on real-valued functions satisfy, at the same two fixed inputs,
\[
 \E_\pi z_w(x_\pm)=0,\qquad
 \E_\pi z_w(x_\pm)^2\le C_0,\qquad
 \E_\pi|z_w(x_+)-z_w(x_-)|^2\le K r^H,
 \qquad 0<r<1.
\]
Use the raw Gaussian likelihood
$\exp[-\sum_i(z_w(X_i)-Y_i)^2/T]$ and its raw posterior mean.
For the noiseless labels $Y=\pm a$ above, every realized sample satisfies
\[
 |\bar z_S(x_+)-\bar z_S(x_-)|
 \le \exp\!\left(\frac{n(C_0+a^2)}{T}\right)\sqrt K\,r^{H/2}.
\]
Consequently, for sufficiently large $H$, risk is at least $a^2/4$
whenever
\[
 \frac nT\le
 \frac{H\log(1/r)}{4(C_0+a^2)}.
\]
There is no bounded-output link and no restriction excluding native biases.
\end{proposition}
\begin{proof}
The expected squared training loss at each observed point is at most
$C_0+a^2$. Jensen's inequality therefore gives
$Z_S\ge\exp[-n(C_0+a^2)/T]$.
The likelihood is at most one, so the raw posterior has finite first
moments and its mean increment is bounded by this inverse normalizer
times the prior mean absolute increment. Cauchy--Schwarz gives the stated
inequality. Under the displayed sample condition it is at most
$\sqrt K r^{H/4}$, eventually at most $a$; the same two-point risk identity
used in Lemma~\ref{lem:increment-learning-lower} finishes the proof.
\end{proof}

\begin{corollary}[Continuous cube and sphere laws]
\label{cor:continuous-raw-depth-lower}
Fix a probability input law $D$ and a bounded noiseless target $g$,
with $|g|\le a$ and $\operatorname{Var}_D(g)=v_g>0$.
Suppose the raw prior is centered pointwise, has
$\sup_x\E_\pi z_w(x)^2\le C_0$, and satisfies
\[
 \E_\pi\operatorname{Var}_D(z_w)\le K r^H,\qquad 0<r<1.
\]
For every realized sample, its raw posterior mean obeys
\[
 \left\|\bar z_S-\E_D\bar z_S\right\|_{L^2(D)}
 \le e^{n(C_0+a^2)/T}\sqrt K\,r^{H/2}.
\]
Hence $R(\bar z_S)\ge v_g/4$ for all sufficiently large $H$ whenever
$n/T\le H\log(1/r)/[4(C_0+a^2)]$.

The independent and neuron-gated tanh priors in the examples below
satisfy the variance premise on every fixed compact input law,
in the same contractive variance regimes. In particular one may use
uniform measure on $[-1,1]^d$ or on the unit sphere, at any fixed
$d\ge2$, and the exact width-one, one-hidden-layer tanh target
$g(x)=a\tanh(x_1)$. Its variance is a positive constant.
The linear target $g(x)=a x_1$, implemented by two ReLU hidden units,
also gives the conclusion.
\end{corollary}
\begin{proof}
Let $\mathcal C f=f-\E_D f$ be the centering operator on $L^2(D)$.
The same Jensen lower bound for the normalizer holds because all
observed labels have absolute value at most $a$.
Minkowski's inequality and Cauchy--Schwarz give
\[
 \|\mathcal C\bar z_S\|_2
 \le\E_{\Pi_S}\|\mathcal C z_w\|_2
 \le Z_S^{-1}\E_\pi\sqrt{\operatorname{Var}_D(z_w)}
 \le e^{n(C_0+a^2)/T}\sqrt K\,r^{H/2}.
\]
The moment assumptions justify the integrals and centering operation.
Risk decomposes into the squared difference of means and a centered
squared distance. If the last bound is at most $\sqrt{v_g}/2$,
the reverse triangle inequality gives $R(\bar z_S)\ge v_g/4$.
The sample condition makes that bound at most $\sqrt K r^{H/4}$.

For the specified priors the same first-layer calculation for two
arbitrary inputs gives a factor bounded by
$C\|x-x'\|_2^2$, followed by the stated contraction in each hidden
layer and a fixed final-readout factor. Consequently
\[
 \E_\pi\operatorname{Var}_D(z_w)
 =\tfrac12\E_{X,X'\sim D}\E_\pi|z_w(X)-z_w(X')|^2
 \le K_D r^H
\]
on any fixed compact domain, after changing the constant $K_D$.
The chosen targets are nonconstant under the uniform cube and sphere
laws, which proves their strictly positive variance.
\end{proof}

\paragraph{Examples with bounded tanh hidden features.}
For the independent variance-gamma fan-in prior, take
$W_F=L_{1/F}$ with
$L_t=\Gamma(ct,1)-\Gamma(ct,1)$, and let native internal and output biases
have the time-one law. Then
$\operatorname{Var}(W_F)=2c/F$,
$\E z_w(x)^2\le4c$, and
\[
 \E|z_w(x_+)-z_w(x_-)|^2\le4(2c)^{H+1}.
\]
Thus every fixed $c<1/2$ supplies the raw lower bound with native biases,
uniformly over all learner widths. The same conclusion holds with a
different independent hidden-weight law of centered fan-in variance
$v/F$, $v<1$, and any independent centered readout law of total variance
bounded uniformly in width.

For the normalized Gaussian readout, let
\[
 z=G_0+\sum_j\sqrt{P_j}G_jh_j^H,
 \qquad P\sim\operatorname{Dirichlet}(\kappa/N,\ldots,\kappa/N).
\]
Its raw second moment is at most two. Suppose the hidden edge factors
have variance $\mu_2$, the hidden gamma gates have shape $\kappa/N$,
and set $r=\kappa(1+\kappa)\mu_2<1$.
The first-layer difference is at most $4\mu_2$ and each subsequent hidden
layer contracts it by at most $r$. The normalized readout contributes
a factor one, since its proportions are independent of their incoming
features and $\E P_j=1/N$. Therefore
\[
 \E|z_w(x_+)-z_w(x_-)|^2
 \le4\mu_2r^{H-1}.
\]
This verifies all hypotheses of Proposition~\ref{prop:raw-increment-depth-lower};
one may take $K=4\mu_2/r$.
The unnormalized Gaussian neuron-gate readout also has uniformly bounded
raw second moment and satisfies the same criterion, with its final
variance factor included.

These are fixed, simple target problems: for tanh, the two labels are
implemented by the width-one, one-hidden-layer function
$g(x)=a\tanh(x_1)/\tanh(1)$.
If a chosen likelihood schedule obeys $T_n\ge c_T n^\theta$ with
$0\le\theta<1$, the necessary sample count for vanishing risk along
increasing depths is
\[
 n=\Omega\!\left(H^{1/(1-\theta)}\right).
\]
In particular $T_n$ proportional to $\sqrt n$ incurs a quadratic-depth
obstruction for that schedule, while $T_n$ proportional to $n^{2/3}$
incurs a cubic-depth obstruction. This is not a lower bound over all
temperature choices or all prior intensities.

The same conclusion covers the power-readout family
$z=G_0+\sum_jP_j^{1-1/p}G_jh_j^H$, with factor density proportional
to $e^{-|G|^p}$ and fixed $p>1$.
Indeed, for $a_N=\kappa/N$ and $q=2-2/p>0$,
\[
 \E\sum_jP_j^q=
 N\frac{\Gamma(a_N+q)\Gamma(\kappa)}
         {\Gamma(a_N)\Gamma(\kappa+q)}
 \le C_{\kappa,p}.
\]
The raw second moment and final increment gain are therefore uniformly
bounded in width. With the same contractive hidden-gate regime, its
prescribed schedule $T_n\asymp n^{2/(p+2)}$ needs
$n=\Omega(H^{1+2/p})$ for small risk on this fixed shallow target.
Constants may depend on $p$; this is again a statement about that
particular schedule.

At fixed likelihood variance the main lower bound is $\Omega(H)$, matching up to logarithms the additive learner-depth upper bounds for shallow teachers. It requires no increasing input dimension, label noise, or complicated target. It does not preclude a depth-dependent likelihood temperature that tends to zero.
